\documentclass[11pt]{article}
\usepackage[a4paper,margin=2.8cm]{geometry}
\usepackage[T1]{fontenc}
\usepackage[utf8]{inputenc}
\usepackage[spanish,es-nodecimaldot]{babel}
\usepackage{lmodern,microtype}
\usepackage{amsmath,amssymb,amsthm,mathtools}
\usepackage{booktabs,array}
\usepackage[table]{xcolor}
\usepackage[most]{tcolorbox}
\usepackage{hyperref}
\usepackage{fancyhdr}

% Paleta institucional UP
\definecolor{upgold}{HTML}{B9975B}
\definecolor{upblue}{HTML}{1F3864}
\definecolor{upwine}{HTML}{6E1423}
\definecolor{upgreen}{HTML}{3B5B40}

\hypersetup{colorlinks=true,linkcolor=upblue,urlcolor=upblue,citecolor=upblue}
\setlength{\parindent}{0pt}\setlength{\parskip}{5pt}

\newtheorem{teorema}{Teorema}[section]
\newtheorem{proposicion}[teorema]{Proposición}
\theoremstyle{definition}
\newtheorem{definicion}[teorema]{Definición}
\newtheorem{ejemplo}[teorema]{Ejemplo}

\newcommand{\R}{\mathbb R}
\newcommand{\D}{\frac{d}{dx}}

% Caja de "Lectura financiera" para conectar cada método con su uso en finanzas cuantitativas
\newtcolorbox{finanzas}[1][]{%
  colback=upgold!8, colframe=upgold!70!black,
  boxrule=0.6pt, arc=1.5pt, left=6pt, right=6pt, top=4pt, bottom=4pt,
  fonttitle=\bfseries\color{upblue}, title={Lectura financiera},#1
}

\setlength{\headheight}{14pt}
\pagestyle{fancy}\fancyhf{}
\lhead{\small\color{upblue} EDO de orden superior en finanzas cuantitativas}
\rhead{\small\thepage}
\renewcommand{\headrulewidth}{0.4pt}
\renewcommand{\headrule}{\color{upgold}\hrule width\headwidth height\headrulewidth}

\title{\textsc{\color{upblue}Ecuaciones Diferenciales Aplicadas}\\[4pt]
\Large Ecuaciones diferenciales de orden superior en finanzas cuantitativas\\[5pt]
\large Modelos de bonos, tasas, ciclos económicos y opciones perpetuas}
\author{Julio César Galindo López\\[2pt]\small Facultad de Ciencias Económicas y Empresariales, Universidad Panamericana}
\date{}

\begin{document}
\maketitle
\begin{center}\emph{Métodos, ejemplos resueltos y aplicaciones a finanzas cuantitativas.}\end{center}

\section*{Introducción}
Muchos modelos centrales de finanzas cuantitativas son, en el fondo, ecuaciones diferenciales
lineales de orden superior. La dinámica de precios de bonos cupón cero bajo tasas
deterministas, los modelos de ciclo económico tipo multiplicador--acelerador, la
reducción de la ecuación de Black--Scholes para reclamos perpetuos, y los ajustes
de precio hacia el equilibrio con corrección de error, comparten la misma pregunta:
\emph{¿qué estructura tiene la ecuación y qué método permite encontrar sus soluciones?}

Estas notas están organizadas alrededor de esa pregunta. Se introducen las herramientas
mínimas y, de inmediato, se aplican a ejemplos financieros de dificultad creciente. La
teoría aparece sólo en la medida necesaria para entender por qué funciona cada
procedimiento, y cada método se cierra con una caja de \textbf{lectura financiera} que
explica dónde aparece en la práctica.

\tableofcontents

\section{La forma general}
Una ecuación de orden $n$ puede escribirse como
\[
F(t,y,y',\ldots,y^{(n)})=0.
\]
En el caso lineal:
\[
a_n(t)y^{(n)}+\cdots+a_1(t)y'+a_0(t)y=g(t).
\]
Si $g=0$ es homogénea; si $g\ne0$, no homogénea.

\begin{ejemplo}
Clasifica los siguientes modelos, donde $P(t)$ es el precio de un activo y $t$ es tiempo:
\[
P''-3P'+2P=0,\qquad t^2P''+tP'-4P=t^3,\qquad P''+\sin P=0.
\]
Las dos primeras son lineales; la primera es homogénea (dinámica libre de precio),
la segunda no homogénea (con un forzamiento exógeno $t^3$).
La tercera es no lineal y aparece típicamente en modelos con fricciones o restricciones
de arbitraje no lineales.
\end{ejemplo}

\begin{finanzas}
En un mercado sin fricciones, muchas cantidades financieras (precio de un bono, valor
presente de un flujo, desviación del PIB respecto de su tendencia) obedecen ecuaciones
lineales: la superposición de estrategias sigue siendo una estrategia válida. La parte
no homogénea $g(t)$ representa siempre un \emph{insumo exógeno}: cupones, dividendos,
gasto público o choques de política monetaria.
\end{finanzas}

\section{Problemas de valores iniciales}
Para orden $n$ se prescriben normalmente
\[
y(t_0)=y_0,\quad y'(t_0)=y_1,\quad\ldots,\quad y^{(n-1)}(t_0)=y_{n-1}.
\]

\begin{teorema}[Existencia y unicidad]
Si los coeficientes de una ecuación lineal son continuos en un intervalo $I$ y el
coeficiente principal no se anula en $I$, el problema de valores iniciales tiene una
única solución en $I$ para datos iniciales dados.
\end{teorema}

\begin{ejemplo}[Capital de un fondo con aportaciones lineales]
Un fondo de inversión tiene un capital $K(t)$ que crece por rendimiento compuesto y
recibe aportaciones que aumentan linealmente en el tiempo. El modelo es
\[
K''-4K=12t,\qquad K(0)=4,\quad K'(0)=1.
\]
La homogénea da $K_h=c_1e^{2t}+c_2e^{-2t}$. Probamos $K_p=At+B$ y obtenemos
$A=-3$, $B=0$. Por tanto
\[
K=c_1e^{2t}+c_2e^{-2t}-3t.
\]
Las condiciones iniciales producen
\[
c_1+c_2=4,\qquad 2c_1-2c_2-3=1,
\]
de donde $c_1=3$, $c_2=1$. Así,
\[
\boxed{K(t)=3e^{2t}+e^{-2t}-3t}.
\]
\end{ejemplo}

\begin{ejemplo}[Ciclo de negocios sin fricción]
En una versión simplificada del modelo multiplicador--acelerador de Samuelson, la
desviación del ingreso $x(t)$ respecto de su nivel de equilibrio, sin amortiguamiento,
satisface
\[
x''+\omega^2x=0,\qquad x(0)=x_0,\quad x'(0)=v_0.
\]
Como $x=C_1\cos(\omega t)+C_2\sin(\omega t)$,
\[
\boxed{x(t)=x_0\cos(\omega t)+\frac{v_0}{\omega}\sin(\omega t)}.
\]
\end{ejemplo}

\begin{finanzas}
Todo problema de valores iniciales corresponde a un modelo donde se conoce el
\emph{estado presente} (precio hoy, velocidad de cambio hoy) y se pide proyectar hacia
adelante: es la estructura natural de cualquier simulación de trayectoria de precio o
de forecasting de una serie macroeconómica a partir de sus condiciones actuales.
\end{finanzas}

\section{Problemas de valores en la frontera}
Se prescriben condiciones en puntos distintos, por ejemplo
\[
y(a)=\alpha,\qquad y(b)=\beta.
\]
Puede haber cero, una o infinitas soluciones.

\begin{ejemplo}[Precio de un bono entre emisión y vencimiento]
Sea el precio de un instrumento
\[
P(t)=c_1t^2+c_2t^4+3.
\]
Si se exige $P(-1)=0$ y $P(1)=4$ (una condición inconsistente con la simetría del
modelo), se obtiene $c_1+c_2=-3$ y $c_1+c_2=1$: \textbf{no existe solución} —el
instrumento no puede satisfacer ambos precios simultáneamente.

Si se exige $P(0)=3$, $P(1)=0$, sólo aparece la condición $c_1+c_2=-3$: hay
\textbf{infinitas soluciones} —el precio en fronteras no identifica de forma única la
trayectoria.

Si se fija el precio de emisión $P(1)=3$ y el precio a un año $P(2)=15$:
\[
c_1+c_2=0,\qquad 4c_1+16c_2=12,
\]
de donde $c_1=-1$, $c_2=1$ y
\[
\boxed{P(t)=t^4-t^2+3}.
\]
\end{ejemplo}

\begin{finanzas}
Los problemas de frontera son el marco natural de la valuación de derivados: en la
ecuación de Black--Scholes se fija el valor terminal (el \emph{payoff} en $T$) y, según
el instrumento, una condición en la frontera del dominio de precios (por ejemplo, el
valor de una opción cuando el subyacente vale cero). Que el problema tenga cero, una o
infinitas soluciones es exactamente la pregunta de si el instrumento está bien
especificado.
\end{finanzas}

\section{El operador diferencial}
Definimos
\[
D=\frac{d}{dt}.
\]
Entonces $D^2y=y''$, $D^3y=y'''$, etc. Por ejemplo,
\[
y''-2y'+5y=0
\quad\Longleftrightarrow\quad
(D^2-2D+5)y=0.
\]

\begin{ejemplo}[Factorización de un modelo de ajuste de precio]
\[
P''-5P'+6P=0
\quad\Longleftrightarrow\quad
(D-2)(D-3)P=0.
\]
Por tanto
\[
\boxed{P=c_1e^{2t}+c_2e^{3t}}.
\]
\end{ejemplo}

\section{Superposición}
\begin{teorema}[Principio de superposición]
Si $y_1,\ldots,y_k$ son soluciones de una ecuación lineal homogénea, entonces
\[
y=c_1y_1+\cdots+c_ky_k
\]
también es solución.
\end{teorema}

\begin{ejemplo}[Combinación de estrategias]
Para el modelo de precio
\[
P''+2P'+P=0
\]
son soluciones $e^{-t}$ y $te^{-t}$. Luego cualquier combinación lineal de ambas
trayectorias también es una trayectoria admisible:
\[
\boxed{P=c_1e^{-t}+c_2te^{-t}}.
\]
\end{ejemplo}

\begin{finanzas}
La superposición es la base matemática de la construcción de un portafolio: si dos
trayectorias de valor son consistentes con la misma ecuación de valuación lineal,
cualquier combinación —cualquier ponderación entre dos activos o dos estrategias
replicantes— también lo es. Esto es lo que permite construir portafolios réplica en
mercados completos.
\end{finanzas}

\section{Independencia lineal y Wronskiano}
Las funciones $f_1,\ldots,f_n$ son linealmente independientes si
\[
c_1f_1+\cdots+c_nf_n=0
\]
implica $c_1=\cdots=c_n=0$.

Para dos funciones:
\[
W(f_1,f_2)=
\begin{vmatrix}f_1&f_2\\f_1'&f_2'\end{vmatrix}
=f_1f_2'-f_1'f_2.
\]
Si $W(t_0)\ne0$, son independientes.

\begin{ejemplo}[Dos factores de descuento independientes]
Para $f_1=e^{2t}$ y $f_2=e^{-t}$ (dos factores de descuento con tasas distintas),
\[
W=-3e^t\ne0.
\]
Luego son linealmente independientes en $\R$: ninguno es múltiplo redundante del otro.
\end{ejemplo}

\begin{ejemplo}[Factores redundantes]
\[
f_1=4t^3,\quad f_2=2t^2,\quad f_3=8t^3+12t^2.
\]
Como
\[
f_3=2f_1+6f_2,
\]
el conjunto es dependiente: $f_3$ no aporta información nueva a la base de riesgo.
\end{ejemplo}

\begin{finanzas}
En la descomposición de la curva de tasas o de un modelo factorial de retornos, el
Wronskiano es la prueba formal de que un conjunto de factores propuesto (por ejemplo,
nivel, pendiente y curvatura de la curva de tasas) es genuinamente independiente y no
contiene una combinación redundante de los otros.
\end{finanzas}

\section{Conjuntos fundamentales}
Un conjunto de $n$ soluciones linealmente independientes de una ecuación homogénea de
orden $n$ es un conjunto fundamental. Si $\{y_1,\ldots,y_n\}$ es fundamental,
\[
\boxed{y=c_1y_1+\cdots+c_ny_n}.
\]

\section{Coeficientes constantes}
Para
\[
a_ny^{(n)}+\cdots+a_1y'+a_0y=0
\]
probamos $y=e^{rt}$. Aparece el polinomio característico
\[
p(r)=a_nr^n+\cdots+a_1r+a_0.
\]
En un modelo de precio o de tasa, las raíces $r$ son \textbf{tasas de crecimiento o
decaimiento propias del sistema}: qué tan rápido el precio regresa a su equilibrio,
crece o oscila sin intervención externa.

\subsection{Raíces reales distintas}
\begin{ejemplo}[Dos componentes de reversión a la media]
\[
P''-5P'+6P=0.
\]
Como $(r-2)(r-3)=0$,
\[
\boxed{P=c_1e^{2t}+c_2e^{3t}}.
\]
\end{ejemplo}

\begin{ejemplo}
\[
y'''-y''-6y'+6y=0.
\]
El polinomio es $(r-1)(r^2-6)$, así que
\[
\boxed{y=c_1e^t+c_2e^{\sqrt6t}+c_3e^{-\sqrt6t}}.
\]
\end{ejemplo}

\begin{finanzas}
Raíces reales distintas describen un instrumento cuyo comportamiento es una mezcla de
crecimientos exponenciales puros —típico de la descomposición espectral de una curva
de tasas o de un modelo de reversión a la media con más de un factor de velocidad de
ajuste.
\end{finanzas}

\subsection{Raíces repetidas}
Si $r$ tiene multiplicidad $m$, aparecen
\[
e^{rt},te^{rt},\ldots,t^{m-1}e^{rt}.
\]

\begin{ejemplo}[Ajuste críticamente amortiguado hacia el equilibrio]
\[
P''-4P'+4P=0
\quad\Longrightarrow\quad
(r-2)^2=0.
\]
Por tanto
\[
\boxed{P=(c_1+c_2t)e^{2t}}.
\]
\end{ejemplo}

\begin{ejemplo}
\[
y'''-3y''+3y'-y=0.
\]
Como el polinomio es $(r-1)^3$,
\[
\boxed{y=e^t(c_1+c_2t+c_3t^2)}.
\]
\end{ejemplo}

\begin{finanzas}
Una raíz repetida corresponde al caso límite entre el ajuste oscilatorio y el
puramente exponencial: el llamado \emph{amortiguamiento crítico}, el punto de
transición más rápido posible sin sobre--corrección de precio hacia el equilibrio, un
caso de referencia en modelos de corrección de error.
\end{finanzas}

\subsection{Raíces complejas}
Si $r=\alpha\pm i\beta$, las soluciones reales son
\[
e^{\alpha t}\cos(\beta t),\qquad
e^{\alpha t}\sin(\beta t).
\]

\begin{ejemplo}[Oscilación amortiguada del precio alrededor del equilibrio]
\[
P''-2P'+5P=0.
\]
Las raíces son $1\pm2i$. Entonces
\[
\boxed{P=e^t(c_1\cos2t+c_2\sin2t)}.
\]
\end{ejemplo}

\begin{ejemplo}[PVI]
Resolver
\[
y''-2y'+5y=0,\qquad y(0)=1,\quad y'(0)=0.
\]
De $y(0)=1$ resulta $c_1=1$ y de $y'(0)=0$ resulta $c_2=-1/2$. Así,
\[
\boxed{y=e^t\left(\cos2t-\frac12\sin2t\right)}.
\]
\end{ejemplo}

\begin{finanzas}
Las raíces complejas son la firma matemática de los \textbf{ciclos económicos}: el
sistema oscila alrededor de su equilibrio con una frecuencia $\beta$ (periodicidad del
ciclo) y una envolvente $e^{\alpha t}$ que indica si el ciclo se amortigua ($\alpha<0$,
mercado estable), se mantiene ($\alpha=0$, ciclo persistente) o se amplifica
($\alpha>0$, inestabilidad creciente).
\end{finanzas}

\section{Ecuaciones no homogéneas}
Para
\[
\mathcal L(y)=g(t)
\]
la estructura es
\[
\boxed{y=y_h+y_p}.
\]
La parte $y_h$ es la dinámica libre del sistema financiero; $y_p$ es la respuesta al
insumo exógeno $g(t)$: un choque de política, un flujo de dividendos, un estímulo
fiscal.

\subsection{Coeficientes indeterminados}
\begin{ejemplo}[Tendencia de dividendos creciente linealmente]
\[
y''-3y'+2y=t.
\]
La homogénea es $c_1e^t+c_2e^{2t}$. Con $y_p=At+B$:
\[
2At+2B-3A=t,
\]
de donde $A=1/2$, $B=3/4$. Por tanto
\[
\boxed{y=c_1e^t+c_2e^{2t}+\frac{t}{2}+\frac34}.
\]
\end{ejemplo}

\begin{ejemplo}[Estímulo exógeno con tasa propia distinta]
\[
y''-3y'+2y=e^{3t}.
\]
Con $y_p=Ae^{3t}$:
\[
(9-9+2)A=1,
\]
por lo que $A=1/2$. Así,
\[
\boxed{y=c_1e^t+c_2e^{2t}+\frac12e^{3t}}.
\]
\end{ejemplo}

\begin{ejemplo}[Resonancia: burbuja de precio]
\[
y''-3y'+2y=e^t.
\]
Como $e^t$ ya aparece en $y_h$, probamos $y_p=Ate^t$. Se obtiene $A=-1$, luego
\[
\boxed{y=c_1e^t+c_2e^{2t}-te^t}.
\]
\end{ejemplo}

\begin{ejemplo}[Demanda estacional]
\[
y''+y=\cos2t.
\]
Tomando $y_p=A\cos2t+B\sin2t$:
\[
-3A\cos2t-3B\sin2t=\cos2t.
\]
Así $A=-1/3$, $B=0$, y
\[
\boxed{y=c_1\cos t+c_2\sin t-\frac13\cos2t}.
\]
\end{ejemplo}

\begin{finanzas}
El caso de resonancia —cuando el insumo exógeno crece exactamente a la tasa natural
del sistema— es el análogo matemático de una \textbf{burbuja o una espiral
inflacionaria}: un estímulo que refuerza en lugar de contrarrestar la dinámica propia
del sistema produce una respuesta $y_p$ que crece más rápido que cualquiera de los dos
factores por separado (aparece el término extra $t\cdot e^{rt}$), en vez de estabilizarse.
\end{finanzas}

\section{Cauchy--Euler: la ecuación detrás de Black--Scholes}
Para
\[
at^2y''+bty'+cy=0,\qquad t>0,
\]
se prueba $y=t^r$. La ecuación indicial es
\[
ar(r-1)+br+c=0.
\]

\begin{ejemplo}
\[
S^2y''-3Sy'+4y=0.
\]
La ecuación indicial es $(r-2)^2=0$, así que
\[
\boxed{y=S^2(c_1+c_2\ln S)}.
\]
\end{ejemplo}

\begin{ejemplo}
\[
S^2y''+Sy'+y=0.
\]
La ecuación indicial es $r^2+1=0$. Por tanto
\[
\boxed{y=c_1\cos(\ln S)+c_2\sin(\ln S)}.
\]
\end{ejemplo}

\begin{finanzas}
Esta sección no es una analogía: es la ecuación real detrás de la valuación de
\textbf{reclamos perpetuos} (opciones americanas perpetuas, valuación de proyectos con
opción real de espera indefinida). Para un derivado independiente del tiempo
$V(S)$ sobre un subyacente con precio $S$, la ecuación de Black--Scholes
\[
\tfrac12\sigma^2S^2V''+(r_f-\delta)SV'-r_fV=0
\]
es exactamente de tipo Cauchy--Euler en la variable $S$. Al sustituir $V=S^r$ se
obtiene una ecuación indicial cuadrática en $r$ cuyas raíces determinan el exponente
de la solución de \emph{barrera libre óptima} en el ejercicio de una opción americana
perpetua. Es el mismo procedimiento que en los ejemplos anteriores, con $t$ reemplazado
por el precio del subyacente $S$.
\end{finanzas}

\section{Cómo elegir el método}
\begin{enumerate}
\item ¿Es homogénea? Resuelve primero la dinámica libre (homogénea).
\item ¿Los coeficientes son constantes? Usa el polinomio característico.
\item ¿Aparecen $S^2y''$, $Sy'$ y $y$ (variable en el subyacente)? Prueba Cauchy--Euler.
\item ¿El insumo exógeno es polinómico, exponencial o trigonométrico? Prueba
coeficientes indeterminados.
\item ¿Hay resonancia (riesgo de burbuja)? Multiplica el ensayo particular por una
potencia de $t$.
\item ¿Hay datos iniciales (precio y velocidad hoy)? Determina las constantes al final.
\item ¿Hay datos en dos extremos (emisión y vencimiento)? Resuelve el sistema de
condiciones de frontera.
\end{enumerate}

\begin{center}
\begin{tabular}{>{\raggedright}p{0.32\textwidth}p{0.55\textwidth}}
\toprule
\rowcolor{upgold!15}
\textbf{Lo que ves} & \textbf{Qué probar primero}\\
\midrule
Coeficientes constantes & ecuación característica $p(r)=0$\\
Raíz $r$ de multiplicidad $m$ & $e^{rt},te^{rt},\ldots,t^{m-1}e^{rt}$\\
Raíces $\alpha\pm i\beta$ & ciclo amortiguado $e^{\alpha t}\cos\beta t,\ e^{\alpha t}\sin\beta t$\\
$aS^2y''+bSy'+cy$ (var.\ en el subyacente) & sustitución $y=S^r$ (Cauchy--Euler / Black--Scholes)\\
Insumo polinomial en $t$ & ensayo polinomial\\
Insumo exponencial $e^{at}$ & ensayo $Ae^{at}$\\
Insumo estacional $\cos bt,\sin bt$ & ensayo $A\cos bt+B\sin bt$\\
Resonancia (burbuja) & multiplicar el ensayo por $t^s$\\
Datos en un punto & problema de valores iniciales\\
Datos en dos puntos & problema de frontera (emisión/vencimiento)\\
\bottomrule
\end{tabular}
\end{center}

\section{Un ejemplo completo: reinversión de cupones en resonancia}

Consideremos el precio de un instrumento con reinversión continua de cupones sujeta a
un choque de tasa a la misma frecuencia natural del sistema:
\[
y''-4y'+13y=e^{2t}\cos(3t).
\]

\textbf{Paso 1. Dinámica libre.}
\[
r^2-4r+13=(r-2)^2+9,
\]
por lo que
\[
y_h=e^{2t}(c_1\cos3t+c_2\sin3t).
\]

\textbf{Paso 2. Respuesta al choque.}
El insumo exógeno tiene la misma frecuencia y el mismo factor exponencial que la
dinámica libre. Hay resonancia. El ensayo correcto es
\[
y_p=te^{2t}(A\cos3t+B\sin3t).
\]
Al sustituir se obtiene
\[
A=0,\qquad B=\frac16.
\]
Por tanto
\[
\boxed{
y=e^{2t}(c_1\cos3t+c_2\sin3t)
+\frac{t}{6}e^{2t}\sin3t}.
\]

\begin{finanzas}
El término $t\cdot e^{2t}\sin3t$ crece más rápido que cualquier componente de la
dinámica libre por sí sola: un choque perfectamente sincronizado con la frecuencia
natural del sistema amplifica la trayectoria en lugar de sólo desplazarla —la firma
matemática de un episodio de inestabilidad autoalimentada.
\end{finanzas}

Este ejemplo resume el procedimiento: resolver la dinámica libre, identificar el tipo
de insumo exógeno, detectar resonancia, construir $y_p$ y comprobar.

\section{Errores frecuentes}

\begin{enumerate}
\item \textbf{Raíz repetida.} Si $r$ tiene multiplicidad $m$, aparecen
$e^{rt},te^{rt},\ldots,t^{m-1}e^{rt}$.

\item \textbf{Raíces complejas.} Para $r=\alpha\pm i\beta$, una base real es
\[
e^{\alpha t}\cos\beta t,\qquad e^{\alpha t}\sin\beta t.
\]

\item \textbf{Resonancia.} Si el ensayo particular pertenece a la dinámica libre, hay
que multiplicarlo por una potencia adecuada de $t$ —de lo contrario se subestima el
riesgo de una respuesta amplificada (burbuja).

\item \textbf{Constantes prematuras.} Conviene obtener primero la solución general y
aplicar después las condiciones iniciales o de frontera.

\item \textbf{Sustitución automática.} $y=S^r$ es especialmente natural para
Cauchy--Euler (variable en el precio del subyacente); no es un método universal para
ecuaciones en el tiempo.

\item \textbf{No comprobar.} Sustituir la solución final en la ecuación original es una
verificación rápida y muy eficaz antes de reportar un resultado de valuación.
\end{enumerate}

\section{Mini--colección de ejemplos}

\begin{ejemplo}[Ciclo puro sin amortiguamiento]
\[
y''+9y=0
\quad\Longrightarrow\quad
\boxed{y=c_1\cos3t+c_2\sin3t}.
\]
\end{ejemplo}

\begin{ejemplo}[Ajuste crítico]
\[
y''-8y'+16y=0
\quad\Longrightarrow\quad
\boxed{y=(c_1+c_2t)e^{4t}}.
\]
\end{ejemplo}

\begin{ejemplo}[Ciclo amortiguado]
\[
y''+4y'+13y=0
\quad\Longrightarrow\quad
\boxed{y=e^{-2t}(c_1\cos3t+c_2\sin3t)}.
\]
\end{ejemplo}

\begin{ejemplo}[Estímulo con tasa distinta a la propia]
\[
y''-y=e^{2t}
\quad\Longrightarrow\quad
\boxed{y=c_1e^t+c_2e^{-t}+\frac13e^{2t}}.
\]
\end{ejemplo}

\begin{ejemplo}[Cauchy--Euler en el subyacente]
\[
S^2y''+Sy'-4y=0
\quad\Longrightarrow\quad
r^2-4=0,
\]
por lo que
\[
\boxed{y=c_1S^2+c_2S^{-2}}.
\]
\end{ejemplo}

\section{Ejercicios}
\begin{enumerate}
\item Un fondo tiene dinámica libre $P''-7P'+12P=0$. Resolverla.
\item Resolver $P''+6P'+9P=0$ (ajuste crítico de precio).
\item Un ciclo económico satisface $x''+4x=0$ con $x(0)=2$, $x'(0)=-1$. Encontrar $x(t)$.
\item Resolver $y''-2y'+5y=0$ con $y(0)=0$, $y'(0)=1$ (oscilación amortiguada de precio).
\item Un fondo recibe aportaciones cuadráticas: resolver $y''-y=t^2$.
\item Demanda estacional: resolver $y''+4y=8\cos2t$ e identificar si hay resonancia.
\item Resolver $y''-4y'+4y=e^{2t}$ (¿hay resonancia?).
\item Reclamo perpetuo tipo Cauchy--Euler: resolver $S^2y''-Sy'-3y=0$ para $S>0$.
\item Determinar si $\{1,t,t^2\}$ —una base polinomial de la curva de tasas— es
linealmente independiente.
\item Determinar si $\{e^t,te^t,e^{2t}\}$ es linealmente independiente.
\item Hallar el Wronskiano de $\{\cos t,\sin t\}$ (dos componentes de un ciclo).
\item Bono con precio de emisión y vencimiento fijos: resolver $y''-5y'+6y=0$ con
$y(0)=1$, $y(1)=0$.
\item Estudiar el problema de frontera $y''+\pi^2y=0$, $y(0)=0$, $y(1)=0$: ¿cuántas
soluciones tiene y qué significa para la especificación del instrumento?
\item Resolver $y'''-6y''+11y'-6y=0$ (tres factores de descuento).
\item Para $y''+2y'+y=0$ (ajuste crítico de precio), construir una base de soluciones
y verificar su independencia con el Wronskiano.
\end{enumerate}

\section*{Resumen de fórmulas}

Para una ecuación lineal homogénea con coeficientes constantes,
\[
a_ny^{(n)}+\cdots+a_1y'+a_0y=0,
\]
se busca $y=e^{rt}$ y se resuelve $p(r)=0$.

\[
\begin{array}{c|c}
\text{raíces} & \text{soluciones}\\
\hline
r_1,\ldots,r_n\text{ distintas} & e^{r_1t},\ldots,e^{r_nt}\\
r\text{ de multiplicidad }m & e^{rt},te^{rt},\ldots,t^{m-1}e^{rt}\\
\alpha\pm i\beta & e^{\alpha t}\cos\beta t,\ e^{\alpha t}\sin\beta t
\end{array}
\]

Para una ecuación no homogénea,
\[
\boxed{y=y_h+y_p}.
\]
La dificultad principal está en elegir correctamente $y_p$ y detectar resonancia
(riesgo de amplificación tipo burbuja).

\[
\boxed{
\text{clasificar}
\longrightarrow
\text{resolver la dinámica libre}
\longrightarrow
\text{buscar la respuesta al insumo exógeno}
\longrightarrow
\text{aplicar condiciones}
\longrightarrow
\text{comprobar}
}
\]

\vspace{8pt}
{\small\color{upblue}\emph{Nota:} las ecuaciones de Cauchy--Euler de la Sección~10
son la base algebraica de la ecuación de Black--Scholes para reclamos independientes
del tiempo; los demás modelos (ciclo económico, reversión a la media, resonancia como
burbuja) son analogías pedagógicas estándar en macro-finanzas para dar intuición
económica a cada tipo de raíz.}

\end{document}
